\documentclass[11pt]{article}
\usepackage[margin=1in]{geometry}
\usepackage{amsmath,amssymb,booktabs}
\usepackage[T1]{fontenc}
\usepackage{hyperref}
\newcommand{\pending}[1]{\par\noindent\textbf{Author action:} #1\par}
\begin{document}
\title{Experimental calibration and validation: manuscript draft}
\author{}\date{}\maketitle
\noindent\textit{Draft based on the archived Run 6 pilot and Run 9/10 transfer summaries. Author-action paragraphs and the evidence appendix are editorial notes, not manuscript text. Equation numbering should be adapted to the manuscript. No new calibration or sensitivity experiment was performed for this draft.}

\section*{Experimental calibration and validation design}
Run 6 is used for calibration, and Runs 9 and 10 are used for validation with the calibrated deformation parameters held fixed. The comparisons assess final sheet geometry for different prescribed zigzag geometries and cycle counts.

\begin{center}
\begin{tabular}{lll}\toprule
Evidence set & Role & Cases \\\midrule
Existing pilot & Parameter calibration & Run 6 \\
Existing pilot & Frozen-parameter validation & Runs 9, 10 \\\bottomrule
\end{tabular}
\end{center}
\pending{The researcher-selected allocation is Run 6 calibration and Runs 9/10 validation. Frozen deformation parameters are verified, but audit model-development and preprocessing decisions before claiming independently held-out validation. Confirm experimental dimensions and fixture identities against laboratory records. Run 7 is omitted from presentation, not erased from the experimental record; state the selection rationale. Prospective new-gripper experiments are outside this draft's reported scope.}

\setcounter{section}{3}
\section{Model calibration}
\subsection{Parameterization and fixed model settings}
Calibration identifies equivalent deformation inputs rather than independently measured constitutive constants. The pilot used four dimensionless parameters,
\begin{equation}
\boldsymbol\theta=(g_t,g_b,o_t,o_b),\qquad
 g_{\parallel,\ell}=g_\ell(1+o_\ell),\quad
 g_{\perp,\ell}=g_\ell(1-o_\ell),\quad \ell\in\{t,b\}.
\end{equation}
These inputs prescribe deformation along and transverse to the local toolpath before applying the overlap and history rules. Material properties, track width, history settings, boundary conditions, mesh resolution, and equilibrium tolerances are fixed during calibration. One parameter set is calibrated using Run 6 and retained without refitting for validation against Runs 9 and 10.

All cases use a $200\times300\times0.77$ mm sheet model with a uniform effective track width of $2$ mm; zigzag geometry and cycle count vary between cases. A structured mesh with $2$ mm grid spacing comprises 15,251 vertices and 30,000 triangular elements. Equilibrium is computed using an HLBFGS-based quasi-Newton phase \cite{liu2010hlbfgs} followed by sparse Newton iterations, with TinyAD-based Hessian assembly \cite{schmidt2022tinyad} and regularization when required. The gradient-norm tolerance is $5\times10^{-14}$ in internal solver units. A single unconstrained equilibrium solve follows accumulation of the prescribed deformation history.

\begin{center}
\begin{tabular}{llll}\toprule
Parameter & Lower bound & Upper bound & Selected value \\\midrule
$g_t$ & $5\times10^{-5}$ & $10^{-2}$ & $2.88351\times10^{-4}$ \\
$g_b$ & $10^{-5}$ & $10^{-3}$ & $1.45880\times10^{-5}$ \\
$o_t$ & $-1$ & $1$ & $-0.398916$ \\
$o_b$ & $-1$ & $1$ & $0.960306$ \\\bottomrule
\end{tabular}
\end{center}
\pending{All three archived states confirm the mesh counts, $\nu=0.33$, numerical Young's modulus $E=1$, thickness $0.77$ mm, disabled fixture constraints and one final-only solve. Do not interpret $E=1$ as a measured physical modulus. Verify the layer thickness split and explain energy normalization in the mechanics section. The mesh size is not a mesh-convergence result. Solver acceptance additionally requires finite values and positive curvature. The ordinary gradient uses the shell-energy implementation; TinyAD supplies Hessian assembly.}

\subsection{Geometry comparison and Bayesian optimization}
The calibration objective measures the discrepancy between the scan and predicted shell using a fixed preprocessed scan for the reported search stage. The measured top surface is approximated as the shell midsurface. A rigid Kabsch transformation aligns the deformed model vertices to their corresponding undeformed reference vertices, and both the preprocessed scan and the predicted surface are leveled and centered on the shared undeformed planar footprint $\{(x_i,y_i,0)\}_{i=1}^{N}$ without fitting the scan to the model prediction. Let $\widehat{\mathbf p}_i$ ($i=1,\dots,N$) denote the aligned scan points, $\widehat T$ the triangulated scan surface on the shared $2\text{ mm}$ grid footprint (retaining triangles with edge lengths $\le 3.5\text{ mm}$ and dual-area weights $a_i$), and $\widehat S(\boldsymbol\theta)$ the aligned predicted surface interpolated at the same shared-footprint points $\widehat{\mathbf q}_i(\boldsymbol\theta)$. The shared-footprint double-sided surface RMS discrepancy is
\begin{align}
d_i(\boldsymbol\theta)&=\min_{\mathbf q\in\widehat S(\boldsymbol\theta)}
\|\widehat{\mathbf p}_i-\mathbf q\|_2,\\
r_i(\boldsymbol\theta)&=\min_{\mathbf p\in\widehat T}
\|\widehat{\mathbf q}_i(\boldsymbol\theta)-\mathbf p\|_2,\\
J(\boldsymbol\theta)&=\left\{\frac12\left[
\frac1N\sum_{i=1}^{N}d_i^2+\sum_{i=1}^{N}a_ir_i^2\right]\right\}^{1/2}.
\end{align}
Here, $N$ is the number of retained scan points on the shared footprint and $a_i=A_i/\sum_{k=1}^{N}A_k$ are normalized barycentric dual-area weights of $\widehat T$. Because both directions measure point-to-triangulated-surface distance on the same shared footprint, identical surfaces satisfy $J(\widehat S,\widehat S)=0$ (eliminating the $\sim 1.0\text{ mm}$ floor that arises when full-sheet $200\times 300\text{ mm}$ boundary vertices are compared against an edge-inset scan point cloud). Distances are reported in millimetres.

The same discrepancy is evaluated for Runs 9 and 10 without refitting the deformation parameters. No multi-case calibration objective is used.

Bayesian optimization uses a Gaussian-process surrogate of $-\log[J/(1\,\mathrm{mm})]$ with a Mat\'ern-$5/2$ covariance and separate length scales. Growth magnitudes are searched in logarithmic coordinates and orthotropy coefficients in linear coordinates, normalized to the search box. The verified final continuation selects batches of four candidates using success-aware expected improvement, evaluated in logarithmic form, with a separate probabilistic model of solver completion and a minimum normalized parameter separation of $0.08$. Failed evaluations do not receive artificial objective values. The selected parameters minimize the valid recorded objectives; surrogate uncertainty is not a confidence interval on physical parameters.

The reported fixed-scan search was warm-started from 80 rescored historical records (70 valid objectives and 10 timeouts), rather than a new initial design. It added 160 evaluations in successive allocations of 40, 40 and 80, yielding 151 valid objectives and nine timeouts. The final continuation used batches of four and stopped after its allocated evaluation count; convergence of the parameter estimates was not a stopping requirement. The lowest valid objective occurred at archived attempt 209. Earlier initialization used a 12-point design comprising one prescribed baseline and 11 maximin Latin-hypercube points (seed 2026); those earlier runs underwent reconciliation and rescoring and must not be treated as an unchanged 12-point initialization of the final search.
\pending{The archived final continuation checkpoint and proposal diagnostics verify the success-aware acquisition above. The archived original configuration and checkpoint establish a 12-point baseline-plus-maximin-LHS design, but the later Stage 2 seed was a reconciled 12-record historical set, not identical to those 12 original design points. Earlier search-stage proposal policies are not fully reconstructed: do not claim one uninterrupted ordinary-EI search following a new space-filling design. The consolidated ledger contains 80 rescored training-only prior records (70 valid, 10 timeouts) and 160 current-search records (151 valid, 9 timeouts). The last 80 records contain 74 valid evaluations and six timeouts; the final recovery stage contains 36 records, 33 valid and three timeouts, terminating at its planned count. These are archived records, not a count of every historical solver invocation or infrastructure retry. Add an appropriate Bayesian-optimization citation; TinyAD is not that reference.}

\subsection{Calibration fit and parameter assessment}
Under the shared-footprint double-sided surface metric and corrected toolpath convention, the Attempt-209 parameter set achieved a symmetric surface RMS objective of $0.626$ mm on Run 6, with a scan-to-model surface RMS of $0.630$ mm, a model-to-scan surface RMS of $0.622$ mm, and a scan-to-model 95th-percentile distance of $1.184$ mm (under the historical uncropped point-to-point reverse metric, the archived objective was $1.182$ mm due to the $2\text{ mm}$ scan edge inset and discrete cell-center offset).

The bottom-layer orthotropy estimate, $o_b=0.960306$, lies close to the upper search bound. The bottom-layer growth is also toward the lower end of its logarithmic search interval. These observations motivate sensitivity and parameter-compensation analyses; they do not demonstrate uniqueness or physical identifiability. A reduced-parameter comparison, for example $o_t=o_b$, would be needed to assess whether four independently adjusted parameters are necessary; no such conclusion is drawn here.

For context, the frozen-parameter transfer summaries under the shared-footprint double-sided surface metric give the following distances. These are single-case geometric comparisons, not uncertainty estimates from repeated experiments.
\begin{center}
\begin{tabular}{lrrrr}\toprule
Case & Combined RMS & Scan-to-model RMS & Reverse surface RMS & Forward $P_{95}$ \\\midrule
Run 6 & 0.626 & 0.630 & 0.622 & 1.184 \\
Run 9 & 0.678 & 0.687 & 0.670 & 1.324 \\
Run 10 & 0.719 & 0.726 & 0.713 & 1.389 \\\bottomrule
\end{tabular}\\[3pt]
\small All entries in mm; both directions use shared-footprint triangulated surfaces (historical uncropped point-cloud combined RMS: $1.182$, $1.250$, and $1.248$ mm).
\end{center}
The recorded Run 9 calculation uses a material-frame path offset of $(0,-28.75)$ mm. Its provenance must be distinguished from parameter fitting when assessing transfer validity. The available summaries confirm reuse of the four Run 6 parameters, but do not alone establish that all modeling and preprocessing choices were independent of these transfer scans.

Spatial residual maps and fixed-location sections complement the distance metrics. A height residual $\Delta z(x,y)=z_{\mathrm{model}}(x,y)-z_{\mathrm{experiment}}(x,y)$ is a signed vertical difference on jointly supported coordinates, not the Euclidean distance entering the objective. In the current Run 6 figure draft, sections at $y=-100,0,100$ mm use piecewise-linear interpolation, without a smoothing fit; unsupported coordinates remain gaps. The overlay uses explicitly labeled $12\times$ vertical exaggeration. Directional shading conveys surface orientation, not error magnitude. ``Analytical'' is the requested figure label for the numerical shell-model prediction, not a claim of a closed-form solution.

\pending{Complete sensitivity and shared-orthotropy comparisons before making parameter-necessity or identifiability claims. Objective-history, residual-map/distribution and fixed-section figures are now available under \path{calibration/paper_figures/manuscript_calibration/}, each with a caption, PDF, PNG and provenance. The history separates 80 rescored prior records (70 valid) from 160 expanded-search/continuation records (151 valid); archived attempt order is not wall-clock completion order. The proposed combined Figure 5 below is not yet assembled.}
\paragraph{Proposed Figure 5 caption.}
Calibration on Run 6. (a) Best successfully evaluated objective against cumulative solver evaluations, with unsuccessful evaluations identified. (b) Signed-height residual map over the measured and model-supported domain, with missing coverage distinguished. (c--e) Experiment and model sections at fixed $y$ positions, accompanied by signed differences. Parameters and search bounds are reported separately in the table. All spatial residuals are distinct from the bidirectional optimization objective. \textit{Editorial note: this figure is proposed, not yet assembled.}

\appendix
\section{Evidence and outstanding checks (not manuscript text)}
The draft was checked against the following local records; it does not claim a full reconstruction of historical executable identities or the whole optimization ledger.
\begin{itemize}
\item \path{run/run6_expanded_bo_continue20_recovery/frozen_config.json}: bounds, fixed settings, case and scoring method.
\item \path{run/run6_expanded_bo_continue20_recovery/bayes_summary.json}: selected parameters, metrics, termination manifest and continuation counts.
\item \path{run/run6_expanded_bo_continue20_recovery/evaluations/attempt_0209/case_00/command.json} and \path{sequence.json} in the same directory: actual solver arguments, disabled fixture and prescribed history.
\item \path{calibration/ewcal/material_frame.py} and \path{calibration/ewcal/geometry.py}: rotation, centering, directed distances and weights.
\item \path{calibration/ewcal/batch_bayes.py}: GP fitting; \path{calibration/analysis/success_aware_bo/model.py} and \path{calibration/analysis/success_aware_bo/separation.py}: success-aware acquisition and separation. \path{run/run6_expanded_bo_continue20_recovery/bayes_state.json}: executed final-continuation proposal diagnostics. Earlier stage-specific acquisition and initial-design reconstruction remain incomplete.
\item \path{run/run9_transfer_expanded/run9_transfer_expanded_summary.json} and \path{run/run10_transfer/run10_transfer_summary.json}: transfer parameters and metrics.
\item \path{docs/calibration/bayesian_calibration/bayesian_calibration.tex}: existing alignment/objective note, retained unchanged.
\end{itemize}
This draft replaces the outdated bounded scan-registration and one-way-objective description for the archived pilot only. Prospective experiments, altered boundary conditions and altered equilibrium schedules require their own documented methods and evidence.
\paragraph{Verified fixed settings for transcription.}
The dimensionless mesh setting is $r=0.01$, giving target spacing $2L_xr=2$ mm for $L_x=0.1$ m. Archived sequences specify (length in mm, angle in degrees, strips, repeats): Run 6 $(130,1,68,29)$; Run 9 $(187.5,1,32,33)$; Run 10 $(175,1.5,24,50)$. Material-frame path offsets are $(0,0)$, $(0,-28.75)$ and $(0,-22.5)$ mm, respectively. The Run 6 frozen history settings are hit-count \texttt{voce\_decay}, $\beta=0.2231435513$, zero floor, face-shared history and plastic-hit counting. Frozen preprocessing settings are $2$ mm spacing, $2$ mm edge inset, $20$ mm outlier floor, $0.5$ mm maximum cell-height MAD and at least 1000 retained points; realized preprocessing still requires an audit.

\begin{thebibliography}{9}
\bibitem{liu2010hlbfgs} Y. Liu, \emph{HLBFGS: A hybrid L-BFGS routine}, software, 2010. \url{https://xueyuhanlang.github.io/software/HLBFGS/}.
\bibitem{schmidt2022tinyad} P. Schmidt, J. Born, D. Bommes, M. Campen, and L. Kobbelt, \emph{TinyAD: Automatic Differentiation in Geometry Processing Made Simple}, \emph{Computer Graphics Forum}, 41(5), 2022. Bibliographic fields verified against the bundled TinyAD citation entry.
\end{thebibliography}
\end{document}
